% Structural-diagram panels formatted as subfigures across the text width.
\begin{subfigure}[t]{0.48\textwidth}
  \centering
  \resizebox{0.94\linewidth}{!}{%
    \begin{tikzpicture}[>=Stealth,
      state/.style={draw,circle,fill=teal!8,minimum size=8mm,inner sep=1pt,
        font=\scriptsize}]
      \node[draw,ellipse,fill=orange!12,inner sep=2pt,font=\scriptsize] (g21-theta) at (0,0) {$\theta$};
      \node[state] (g21-state) at (2.1,0) {$x_k$};
      \node[draw,rectangle,rounded corners,fill=blue!7,inner sep=3pt,font=\scriptsize] (g21-obs) at (4.6,0) {$y_k$};
      \node[draw,ellipse,fill=gray!10,inner sep=2pt,font=\scriptsize] (g21-noise) at (4.6,-1.1) {$v_k$};
      \draw[->,thick] (g21-theta)--(g21-state);
      \draw[->,thick] (g21-state)--(g21-obs);
      \draw[->,thick] (g21-noise)--(g21-obs);
    \end{tikzpicture}%
  }
  \caption{有向依赖网络}
\end{subfigure}\hfill
\begin{subfigure}[t]{0.48\textwidth}
  \centering
  \resizebox{0.94\linewidth}{!}{%
    \begin{tikzpicture}[>=Stealth,
      process/.style={draw,rounded corners=2pt,fill=blue!5,minimum width=17mm,
        minimum height=8mm,align=center,font=\scriptsize},
      decision/.style={draw,diamond,aspect=2,fill=orange!8,inner sep=1pt,
        align=center,font=\scriptsize}]
      \node[process] (g21-input) at (0,0) {输入参数};
      \node[decision] (g21-check) at (2.6,0) {误差\\可接受？};
      \node[process] (g21-output) at (5.2,.75) {输出结果};
      \node[process] (g21-refine) at (5.2,-.75) {调整网格};
      \draw[->,thick] (g21-input)--(g21-check);
      \draw[->,thick] (g21-check.east)--node[above,font=\tiny]{是}(g21-output.west);
      \draw[->,thick] (g21-check.south)--node[below left,font=\tiny]{否}
        ++(0,-.35)-|(g21-refine.west);
      \draw[->,thick] (g21-refine.south) -- ++(0,-.4)
        -|(g21-input.south);
    \end{tikzpicture}%
  }
  \caption{判断流程图}
\end{subfigure}

\par\medskip

\begin{subfigure}[b]{0.48\textwidth}
  \centering
  \resizebox{0.82\linewidth}{!}{%
  \begin{tikzpicture}[>=Stealth]
    \node[draw,circle,fill=teal!12,minimum size=13mm,inner sep=1pt,
      font=\scriptsize] (g21-root) at (0,0) {状态空间};
    \node[draw,circle,fill=blue!9,minimum size=10mm,inner sep=1pt,
      font=\scriptsize] (g21-state-concept) at (0,1.55) {状态};
    \node[draw,circle,fill=orange!13,minimum size=10mm,inner sep=1pt,
      font=\scriptsize] (g21-parameter-concept) at (1.65,0) {参数};
    \node[draw,circle,fill=green!10,minimum size=10mm,inner sep=1pt,
      font=\scriptsize] (g21-observation-concept) at (0,-1.55) {观测};
    \node[draw,circle,fill=red!9,minimum size=10mm,inner sep=1pt,
      font=\scriptsize] (g21-error-concept) at (-1.65,0) {误差};
    \draw[thick] (g21-root)--(g21-state-concept);
    \draw[thick] (g21-root)--(g21-parameter-concept);
    \draw[thick] (g21-root)--(g21-observation-concept);
    \draw[thick] (g21-root)--(g21-error-concept);
  \end{tikzpicture}
  }
  \caption{状态空间概念图}
\end{subfigure}\hfill
\begin{subfigure}[b]{0.48\textwidth}
  \centering
  \resizebox{0.94\linewidth}{!}{%
    \begin{tikzpicture}[>=Stealth]
      \draw[->,thick] (0,0)--(5.8,0) node[right,font=\scriptsize] {$t$};
      \draw[thick] (.7,.17)--(.7,-.17);
      \draw[thick] (2.9,.17)--(2.9,-.17);
      \draw[thick] (5.1,.17)--(5.1,-.17);
      \node[font=\tiny,align=center] at (.7,.5) {初始化};
      \node[font=\tiny,align=center] at (2.9,-.55) {状态更新};
      \node[font=\tiny,align=center] at (5.1,.5) {观测};
      \node[font=\scriptsize] at (.7,-.42) {$t_0$};
      \node[font=\scriptsize] at (2.9,-.42) {$t_1$};
      \node[font=\scriptsize] at (5.1,-.42) {$t_2$};
    \end{tikzpicture}%
  }
  \caption{状态更新事件时间线}
\end{subfigure}
